\chapter{ViLT 模型简介}

ViLT（Vision-and-Language Transformer）是一类将图像和文本在同一 Transformer 架构中联合建模的视觉—语言模型，其实现要点包括：
\begin{itemize}
  \item 文本嵌入：使用词表嵌入、位置嵌入与分段（token type）嵌入对输入问题进行编码；嵌入维度通常为 768。
  \item 视觉嵌入：将图像划分为若干 patch 并通过一个小的卷积投影层（例如 kernel=32, stride=32）将每个 patch 投影为与文本相同维度的视觉 token。
  \item 联合编码器：文本 token 与视觉 token 在同一序列中送入多层 Transformer 编码器（如 12 层、每层包含自注意力与前馈网络），以捕捉跨模态交互。
  \item 池化与分类头：编码后通过池化器得到整体表征，随后接入小型 MLP/LayerNorm/GELU 层并映射到输出词表上（在本实验使用的预训练模型，其分类器最终输出维度约为 3129，用于生成开放式答案的 logits）。
\end{itemize}

上述设计的直接后果是：模型在生成答案时使用的是一个大规模的开放域词表，这有利于开放式问答的表达能力，但在面对固定、有限的标签集合时，会增加后处理映射的难度（需要把自由文本答案可靠地映射到预定义类）。


\chapter{图文多模态问答（VQA）任务实验}

\section{实验目标}
本章报告基于视觉—语言模型的图文多模态问答（VQA）实验，并辅以监督的图像分类微调作为基线比较。具体目标为：
\begin{itemize}
  \item 验证基于预训练视觉—语言模型的 VQA 推理流程并保存推理输出与评估指标；
  \item 设计并评估把模型自由文本答案映射到固定类标签的后处理策略（精确匹配、模糊匹配与同义词映射）；
  \item 对图像做监督微调（ResNet-18），作为强监督基线并对比不同方法的准确率与误判模式；
  \item 分析误判类型并提出可行的改进路径（例如 CLIP 打分、候选重排序或针对性微调）。
\end{itemize}

\section{方法概述}
本实验采取两条并行技术路线：
\begin{enumerate}
  \item VQA 推理：使用现成的视觉—语言模型对图像与问题对进行前向推理，得到自由文本答案并通过后处理映射为类别标签；
  \item 监督分类：以预训练的图像分类网络为初始化，对目标类别做小规模微调，得到一个可直接输出目标标签的基线分类器。
\end{enumerate}

后处理映射的核心思想是：先对预测文本做规范化（小写、去除标点），做精确匹配；若失败则尝试规则同义词替换（如 "plane"→"airplane"、"car"→"automobile" 等），再使用基于相似度的模糊匹配做候选匹配；若仍无法匹配则将该输出标为未对齐以便人工分析。

\section{数据与预处理}
\subsection{数据集说明}
本实验主要采用公开的小型图像分类数据集作为图像来源（本次实验以 CIFAR-10 为例）。数据集说明如下：
\begin{itemize}
  \item 数据规模：数据集中通常包含 60000 张 32$\times$32 彩色图像，划分为训练集与测试集。
  \item 类别与分布：共有 10 个类别（例如 airplane、automobile、bird、cat、deer、dog、frog、horse、ship、truck），各类别在标准划分中近似均衡（测试集中约 1000 张/类）。
  \item 图像特性：图像分辨率较低，物体通常占据较小像素区域且存在遮挡或复杂背景；这对基于 patch 的视觉 token 表征与小物体识别提出挑战，尤其在将图像上采样以适配预训练视觉—语言模型时需格外注意空间信息的丢失或插值失真。
  \item 金标与缺失：实验评估以数据集提供的类别标签作为金标；若某些推理样本缺少金标或无法确定正确类别，则在评估时排除这些样本以避免对准确率产生偏差。
  \item 子集与快速验证：为便于快速迭代和调试，可使用训练/测试子集进行小规模验证，完整评估则需使用标准的测试拆分以保证统计稳健性。
  \item 对 VQA 的影响：由于预训练视觉—语言模型通常对输入图像做 patch 投影（按固定 kernel 与 stride），输入分辨率和缩放策略会直接影响视觉 token 的空间粒度；实验中对图像做了适当的缩放以保证视觉编码具有足够的空间信息，从而提高跨模态交互的效果。
\end{itemize}

\subsection{预处理}
\begin{itemize}
  \item 对用于视觉—语言模型的输入，统一把图像转换为 PIL.Image，并交由模型的处理器做缩放与归一化；
  \item 对于监督分类器，按预训练权重的输入分布进行数据增强与归一化（例如 Resize 到 224、随机水平翻转、ImageNet 均值/方差归一化）；
\end{itemize}

\section{实验设置}
\subsection{VQA 推理}
使用的预训练视觉—语言模型在开放域 VQA 上表现良好，推理环节的要点如下：
\begin{itemize}
  \item 问题模板：采用通用问题 "What is in the image?"（针对性实验可替换为 "What animal is in the image?"）；
  \item 输出诊断：在推理前后打印并检查模型词表与典型预测，以帮助调试映射策略；
  \item 结果记录：对每条样本保存原始预测文本、映射后的类别与原始标签。
\end{itemize}


\subsection{监督微调（ResNet-18）}
监督微调用于得到可比较的有监督基线，训练要点：
\begin{itemize}
  \item 网络结构：以 ResNet-18 初始化，替换最后的全连接层以适配目标类别数；
  \item 超参数（实验示例）：小规模验证配置使用少量样本（如 2000 条训练 / 1000 条测试）、batch=64、epochs=3、学习率=1e-3；输入按预训练权重做 Resize/Normalize。训练后保存模型权重与训练曲线以供分析。 
\end{itemize}

\section{实验结果}

\subsection{三组训练策略准确率对比}
为清晰展示不同训练策略的效果，下面给出三组实验的准确率对比：
\begin{table}[htbp]
\centering
\begin{tabular}{lcc}
\hline
实验 & 准确率 \\
\hline
未训练（初始权重） & 0.6064 \\
小规模微调 & 0.6790 \\
中等规模微调 & 0.8500 \\
\hline
\end{tabular}
\caption{三种训练策略的准确率比较}
\end{table}

\subsection{未训练ViLT每类准确率}
下面给出 ResNet-18 在小规模微调测试拆分上的每类统计：
\begin{table}[htbp]
\centering
\begin{tabular}{lrrr}
\hline
类别 & 总数 & 正确数 & 准确率 \\
\hline
airplane & 103 & 71 & 0.6893203 \\
automobile & 89 & 87 & 0.9775280 \\
bird & 100 & 87 & 0.8700 \\
cat & 103 & 56 & 0.5436893 \\
deer & 90 & 50 & 0.5555555 \\
dog & 86 & 24 & 0.2790697 \\
frog & 112 & 87 & 0.7767857 \\
horse & 102 & 74 & 0.7254901 \\
ship & 106 & 79 & 0.7452830 \\
truck & 109 & 64 & 0.5871559 \\
\hline
\end{tabular}
\caption{ResNet-18 小规模微调在测试拆分上的每类统计}
\end{table}

\section{实验结果分析}
\subsection{总体表现}
本节汇总本次实验的关键数值与总体观察：
\begin{itemize}
  \item ResNet-18 未训练基线在评估拆分上的准确率为 0.606430155210643。
  \item ResNet-18 小规模微调（train[:2000]）在测试拆分上的准确率为 0.6790。
  \item 中等规模微调（train[:10000]）在验证集上的最佳验证准确率为 0.8500。
\end{itemize}

\subsection{每类性能与混淆概况}
从小规模微调的每类统计可以看出明显的类别间差异（表格中已有详细数值）：
\begin{itemize}
  \item `automobile` 准确率最高（≈0.9775），说明该类别的外观特征对模型较友好；
  \item `dog` 准确率最低（≈0.2791），表明该类在小样本情况下具有高内在变异，易与 `cat` 混淆；
  \item `frog`、`ship`、`horse` 等类别处于中等偏上水平（约 0.72–0.78 区间）。
\end{itemize}

\subsection{混淆矩阵揭示的典型误判包括}

\begin{itemize}
  \item `truck` 常被误判为 `automobile` / `bus`（小样本下类别边界模糊）；
  \item `dog` 与 `cat` 的互相混淆反映了姿态、视角与毛色的多样性带来的识别困难；
  \item 未训练的 VQA 输出包含大量与目标标签无关的短语（如 "pizza","keyboard" 等），说明开放域生成与闭集标签对齐存在较大噪声。。
\end{itemize}

\subsection{方法比较与深入分析}
\begin{itemize}
  \item 监督微调在闭集分类任务上表现稳健，可通过更多带标签数据或更长训练获得进一步提升；
  \item VQA（ViLT）适用于开放式问答与富文本解释场景，但直接用于返回闭集标签需要可靠的对齐机制——当前采用的字符串/模糊匹配易受词汇多样性和描述习惯影响；
  \item 中等规模样本带来的验证精度显著提升（从小样本的 best val≈0.728 到 0.85），表明数据量仍是影响泛化能力的关键因素。
\end{itemize}

\chapter{结论}
\section{实验结论}
\begin{itemize}
  \item 监督微调（基于 ResNet-18）在固定标签集的闭集分类任务中是最直接且有效的方式：小规模微调已显著优于未经对齐的模型，中等规模训练带来更大的提升（验证精度可达 0.85）。
  \item 通用 VQA（ViLT）在开放域任务上具有表达与解释优势，但其开放域输出与固定标签集合的对齐能力有限，需要候选重排序或任务特定微调才能稳定匹配闭集标签。 
  \item 类别间差异提醒我们在数据采集与增强上需做有针对性的设计（特别是高内在变异的类别），并可考虑融合跨模态相似度打分以提高映射鲁棒性。
\end{itemize}

\section{心得体会}
本次实验强化了两点经验：一是“方法要与任务匹配”——开放域的生成模型并不天然适合返回固定标签集的高精度需求；二是“数据与对齐比模型更重要”——在有限模型容量下，合适的数据量、质量与对齐策略通常能带来比更复杂模型更显著的收益。同时学习了ViLT的架构与推理流程，理解了视觉—语言模型在实际应用中的优势与局限，为后续的多模态研究打下了基础。
